\documentclass[11pt]{article}

% -------------------- Packages --------------------
\usepackage[margin=1in]{geometry}
\usepackage{hyperref}
\usepackage{enumitem}
\usepackage{titlesec}
\usepackage{fancyhdr}
\usepackage{xcolor}
\usepackage{listings}
\usepackage{courier}

% -------------------- Hyperref Setup --------------------
\hypersetup{
  colorlinks=true,
  linkcolor=blue,
  urlcolor=blue,
  citecolor=blue
}

% -------------------- Header/Footer --------------------
\pagestyle{fancy}
\fancyhf{}
\lhead{MCV lab}
\rhead{Python Coding Standards}
\cfoot{\thepage}

% -------------------- Section Formatting --------------------
% NOTE: Use titlesec only (do NOT also use sectsty), otherwise formatting conflicts.
\titleformat{\section}{\Large\bfseries}{\thesection}{0.8em}{}
\titleformat{\subsection}{\normalsize\bfseries}{\thesubsection}{0.8em}{}

% Optional: slightly reduce vertical spacing around headings (tweak if desired)
\titlespacing*{\section}{0pt}{1.5em}{0.8em}
\titlespacing*{\subsection}{0pt}{1.0em}{0.5em}

% -------------------- Listings Settings --------------------
\lstset{
  basicstyle=\footnotesize\ttfamily,
  breaklines=true,
  frame=single,
  columns=fullflexible,
  keepspaces=true,
  showstringspaces=false,
  tabsize=4,
  xleftmargin=0.02\textwidth,
  xrightmargin=0.02\textwidth
}

% -------------------- Document Metadata --------------------
\title{\textbf{Medical Computer Vision Lab \\ \normalsize{Python Coding Standards}}}
\author{Maintainer: Ashery Mbilinyi}
\date{Version 1.0 : \today}

\begin{document}
\maketitle

\tableofcontents
\newpage

% -------------------- Abstract (new page + included in ToC) --------------------
\addcontentsline{toc}{section}{Abstract}
\begin{abstract}
This document defines the Python coding standards for \textbf{Medical Computer Vision (MCV) Lab}. We follow \textbf{PEP~8} as the primary style guide and emphasize \textbf{readability, reproducibility, and maintainability}. These standards apply to all research code, prototypes intended for reuse, and production-level systems developed in the lab.
\end{abstract}
\newpage

% ============================================================
\section{Scope and Principles}

\subsection{Scope}
These standards apply to:

\begin{itemize}[label={-},leftmargin=2em,labelsep=0.6em,itemsep=0.2em]
  \item Research prototypes that will be shared within the group
  \item Code released with papers (GitHub repositories)
  \item Internal tooling and shared utilities (datasets, evaluation, plotting)
  \item Any code intended to be maintained beyond the original author
\end{itemize}

\subsection{Guiding Principles}
We follow the \textbf{Zen of Python} (\texttt{import this}) as a guiding philosophy:
\begin{itemize}[label={-},leftmargin=2em,labelsep=0.6em,itemsep=0.2em]
  \item \emph{Readability counts}
  \item \emph{Simple is better than complex}
  \item \emph{Explicit is better than implicit}
\end{itemize}

% ============================================================
\section{Code Layout and Formatting (PEP 8)}

Python coding standards are primarily defined by \textbf{PEP~8}, the official style guide. Adhering to these standards ensures code is maintainable and easily understandable by other developers.

\subsection{Indentation}
\begin{itemize}[label={-},leftmargin=2em,labelsep=0.6em,itemsep=0.2em]
  \item Use \textbf{4 spaces} per indentation level.
  \item Do \textbf{not} use tabs.
\end{itemize}

\subsection{Line Length}
\begin{itemize}[label={-},leftmargin=2em,labelsep=0.6em,itemsep=0.2em]
  \item Limit all lines to a maximum of \textbf{79 characters}.
  \item When lines become too long, prefer:
  \begin{itemize}[label={-},leftmargin=2em,labelsep=0.6em,itemsep=0.15em]
    \item breaking long expressions into multiple lines
    \item using parentheses for implicit line continuation
    \item splitting long strings into smaller parts
  \end{itemize}
\end{itemize}

\subsection{Blank Lines}
\begin{itemize}[label={-},leftmargin=2em,labelsep=0.6em,itemsep=0.2em]
  \item Separate top-level functions and class definitions with \textbf{two blank lines}.
  \item Separate method definitions inside a class with \textbf{one blank line}.
\end{itemize}

\subsection{Imports}
\begin{itemize}[label={-},leftmargin=2em,labelsep=0.6em,itemsep=0.2em]
  \item Imports should be on \textbf{separate lines}.
  \item Group imports in the following order:
  \begin{enumerate}[leftmargin=2em,itemsep=0.15em]
    \item Standard library
    \item Third-party libraries
    \item Local application / library-specific modules
  \end{enumerate}
  \item Separate groups using a \textbf{single blank line}.
\end{itemize}

\textbf{Example:}
\begin{lstlisting}[language=Python]
# Standard library
from pathlib import Path
from typing import Optional

# Third-party
import numpy as np
import torch

# Local
from labproject.data import load_dataset
from labproject.models import build_model
\end{lstlisting}

\subsection{Whitespace}
\begin{itemize}[label={-},leftmargin=2em,labelsep=0.6em,itemsep=0.2em]
  \item Avoid extraneous whitespace immediately inside parentheses/brackets/braces.
  \item Avoid whitespace before commas, semicolons, or colons.
\end{itemize}

\textbf{Preferred:}
\begin{lstlisting}[language=Python]
x = (a + b) * c
items = [1, 2, 3]
my_dict = {"a": 1, "b": 2}
\end{lstlisting}

\textbf{Avoid:}
\begin{lstlisting}[language=Python]
x = ( a + b ) * c
items = [1 , 2 , 3]
my_dict = { "a" : 1 , "b" : 2 }
\end{lstlisting}

% ============================================================
\section{Naming Conventions}

\begin{itemize}[label={-},leftmargin=2em,labelsep=0.6em,itemsep=0.2em]
  \item \textbf{Variables and functions:} \texttt{snake\_case}
  \item \textbf{Classes:} \texttt{PascalCase}
  \item \textbf{Constants:} \texttt{UPPER\_CASE\_WITH\_UNDERSCORES}
  \item \textbf{Modules and packages:} short, lowercase names (underscores allowed if needed)
  \item \textbf{Internal/private usage:} prefix with \texttt{\_} (leading underscore)
\end{itemize}

\textbf{Example:}
\begin{lstlisting}[language=Python]
MAX_EPOCHS = 50

class TraumaClassifier(torch.nn.Module):
    def __init__(self, in_channels: int):
        super().__init__()
        self._in_channels = in_channels

    def forward(self, x: torch.Tensor) -> torch.Tensor:
        return x

def train_one_epoch(model, dataloader):
    for batch in dataloader:
        pass
\end{lstlisting}

% ============================================================
\section{Comments and Documentation}

\subsection{Comments}
\begin{itemize}[label={-},leftmargin=2em,labelsep=0.6em,itemsep=0.2em]
  \item Comments should be \textbf{complete sentences}.
  \item Begin comments with \texttt{\#} followed by a \textbf{single space}.
\end{itemize}

\subsection{Inline Comments}
\begin{itemize}[label={-},leftmargin=2em,labelsep=0.6em,itemsep=0.2em]
  \item Use inline comments sparingly.
  \item Separate inline comments by at least \textbf{two spaces}.
\end{itemize}

\textbf{Example:}
\begin{lstlisting}[language=Python]
step += 1  # Increment global step counter
\end{lstlisting}

\subsection{Docstrings (PEP 257)}
\begin{itemize}[label={-},leftmargin=2em,labelsep=0.6em,itemsep=0.2em]
  \item Use \texttt{"""triple double quotes"""}.
  \item Add docstrings to all public modules, classes, functions, and methods.
  \item Docstrings must clearly describe:
  \begin{itemize}[label={-},leftmargin=2em,labelsep=0.6em,itemsep=0.15em]
    \item purpose
    \item input arguments
    \item return values
    \item major assumptions / side effects
  \end{itemize}
\end{itemize}

\textbf{Example:}
\begin{lstlisting}[language=Python]
def load_ct_volume(path: str) -> np.ndarray:
    """Load a CT volume from disk.

    Args:
        path: Path to a NIfTI file.

    Returns:
        The CT volume as a NumPy array.
    """
    return np.load(path)
\end{lstlisting}

% ============================================================
\section{Programming Recommendations (Pythonic Code)}

\subsection{Comparisons}
\begin{itemize}[label={-},leftmargin=2em,labelsep=0.6em,itemsep=0.2em]
  \item Use \texttt{is None} / \texttt{is not None} for \texttt{None} checks.
\end{itemize}

\textbf{Preferred:}
\begin{lstlisting}[language=Python]
if x is None:
    return
\end{lstlisting}

\textbf{Avoid:}
\begin{lstlisting}[language=Python]
if x == None:
    return
\end{lstlisting}

\subsection{Boolean Values}
\begin{itemize}[label={-},leftmargin=2em,labelsep=0.6em,itemsep=0.2em]
  \item Do not compare booleans to \texttt{True} or \texttt{False} with \texttt{==}.
\end{itemize}

\textbf{Preferred:}
\begin{lstlisting}[language=Python]
if greeting:
    print("Hello")
\end{lstlisting}

\textbf{Avoid:}
\begin{lstlisting}[language=Python]
if greeting == True:
    print("Hello")
\end{lstlisting}

\subsection{String Formatting}
\begin{itemize}[label={-},leftmargin=2em,labelsep=0.6em,itemsep=0.2em]
  \item Prefer \textbf{f-strings} for readable formatting.
\end{itemize}

\begin{lstlisting}[language=Python]
name = "Ashery"
score = 0.92
print(f"{name} achieved AUROC={score:.3f}")
\end{lstlisting}

\subsection{Exceptions}
\begin{itemize}[label={-},leftmargin=2em,labelsep=0.6em,itemsep=0.2em]
  \item Avoid bare \texttt{except:}.
  \item Catch specific exceptions.
\end{itemize}

\textbf{Preferred:}
\begin{lstlisting}[language=Python]
try:
    value = int(x)
except ValueError:
    value = 0
\end{lstlisting}

\subsection{Context Managers}
\begin{itemize}[label={-},leftmargin=2em,labelsep=0.6em,itemsep=0.2em]
  \item Use \texttt{with} statements for resource management (files, etc.).
\end{itemize}

\begin{lstlisting}[language=Python]
from pathlib import Path

path = Path("results.txt")
with path.open("w", encoding="utf-8") as f:
    f.write("done\n")
\end{lstlisting}

% ============================================================
\section{Tools to Enforce Standards}

To ensure consistency across the lab, we rely on automated tools.

\subsection{Recommended Tooling}
\begin{itemize}[label={-},leftmargin=2em,labelsep=0.6em,itemsep=0.2em]
  \item \textbf{Formatter:} \texttt{black} (preferred)
  \item \textbf{Linting:} \texttt{flake8} or \texttt{pylint}
  \item \textbf{Import sorting:} \texttt{isort}
  \item \textbf{Type checking:} \texttt{mypy}
\end{itemize}

\subsection{Minimum Project Requirements}
Every shared lab project repository should include:
\begin{itemize}[label={-},leftmargin=2em,labelsep=0.6em,itemsep=0.2em]
  \item \texttt{README.md} with setup + usage
  \item A reproducible environment file (e.g., \texttt{requirements.txt} or \texttt{environment.yml})
  \item Tool configuration (e.g., \texttt{pyproject.toml} for \texttt{black/isort})
\end{itemize}

% ============================================================
\section{References}

\begin{itemize}[label={-},leftmargin=2em,labelsep=0.6em,itemsep=0.2em]
  \item PEP~8 --- Style Guide for Python Code: \url{https://peps.python.org/pep-0008/}
  \item PEP~257 --- Docstring Conventions: \url{https://peps.python.org/pep-0257/}
  \item Zen of Python: \url{https://peps.python.org/pep-0020/}
  \item Black: \url{https://black.readthedocs.io/}
  \item Flake8: \url{https://flake8.pycqa.org/}
  \item isort: \url{https://pycqa.github.io/isort/}
  \item mypy: \url{https://mypy.readthedocs.io/}
\end{itemize}

\end{document}
